\documentclass[12pt]{article}
\usepackage{setspace}   % paquete para interlineado
\usepackage{newtxtext,newtxmath} % Times New Roman

%otra opcion para times roman adaptable a versiones
%\usepackage{mathptmx}

\setstretch{1.5}  
\usepackage{amsmath} % Para usar \text{}
\usepackage[utf8]{inputenc}
\usepackage[spanish]{babel}
\usepackage[letterpaper, left=3cm, right=2.5cm, top=2.5cm, bottom=2.5cm]{geometry}
\usepackage{graphicx}
\usepackage{titlesec}
\usepackage{lipsum}
\usepackage{enumitem} % Para mejorar listas

\begin{document}
	%Para no generar numeracion de paginas
	\pagestyle{empty}
	
	
	\section*{Datos del estudiante}
	\begin{itemize}[label=$\blacksquare$]
		\item Nombre completo del estudiante:  Marlon Ivan Carreto Rivera
		\item Carrera: Ingeniería Civil
		\item Registro académico: 201230088
		\item Carné (DPI): 2733955050909 
		\item Correo electrónico institucional: marlonivan-carretorivera@cunoc.edu.gt
		\item Tipo de práctica que realizó: Práctica en Docencia
		\item Fecha de inicio de prácticas: 12 de Agosto 2025
		\item Horas de práctica: 116 Horas 
		\item Fecha de finalización de la práctica: 23 de Octubre 2025
	
	\end{itemize}
	
	\section*{Datos del supervisor}
	\begin{itemize}[label=$\blacksquare$]
		\item Nombre de la institución: División de Ciencias de la Ingeniería CUNOC-USAC
		\item Nombre del supervisor de práctica: Erick Adolfo Coti Sac
		\item Correo institucional: erickadolfocoti@cunoc.edu.gt
		\item Profesión: Ingeniero Civil
	\end{itemize}
	
	\section*{Datos del curso}
	\begin{itemize}[label=$\blacksquare$]
		\item Nombre del curso: Laboratorio de Matemática Básica 1
		\item Código del curso: 3000
		\item Cinco actividades principales que le fueron asignadas en la práctica: Se planificó, impartió y evaluó temas como:
		\begin{itemize}
			\item Ecuaciones lineales y números complejos.
			\item Geometría básica y resolución de ejercicios.
			\item Primer examen parcial y uso de GeoGebra.
			\item Funciones polinomiales y su graficación.
			\item Aplicaciones de funciones polinomiales.
		\end{itemize}
		
		
		\item Mencionar si existe un manual o documento en el que se trabajo:  Hay dos guias 
		\begin{itemize}
			\item Programa de Cursos Matematica Basica 1
			\item Castillo, M. Á. (s. f.). Ecuaciones lineales. Recursos Matemática en Línea, Facultad de Ingeniería, Universidad de San Carlos de Guatemala. Recuperado el 2 de noviembre de 2025, de https://matematicaenlinea.com/basica-1/unidad-1-ecuaciones-y-desigualdades/1-1-ecuaciones/
			
		\end{itemize}
		
	\end{itemize}
	
\end{document}
